\documentclass[10pt,a4paper]{article}
\usepackage[utf8]{inputenc}
\usepackage[a4paper,left=3cm,right=2cm]{geometry}
\usepackage{amsmath}
\usepackage[thmmarks, amsmath, thref]{ntheorem}
\usepackage{amsfonts}
\usepackage{amssymb}
\usepackage{fancyhdr}
\pagestyle{fancy}
\usepackage{anysize}
\usepackage{graphicx}
\usepackage{float}
\usepackage{hyperref}
\usepackage{multirow}
\usepackage{enumitem}
\newtheorem{theorem}{Theorem}
\theoremsymbol{\ensuremath{\blacksquare}}
\newtheorem*{solution}{Solución:}
\usepackage{listings}
\usepackage{xcolor}

% Margenes y formalidades

\textwidth 6.25in % Margen Derecho y ancho del texto
\textheight 8.5in
\oddsidemargin 0in % Margen Izquierdo
\headheight 0in % Largo del texto
\leftmargin 1in
\parindent 0em % Salto entre la indentación (Sangría)https://www.overleaf.com/4519127669sdpptkhvsbfn
\topmargin -0.5in % Margen Superior
\parskip 0.5ex % Salto entre parrafos
\baselineskip 1pt

%\lhead{\includegraphics[scale=0.2]{logo.png}} % texto izquierda de la cabecera
%\chead{}                                      % texto centro de la cabecera
%\rhead{\includegraphics[scale=0.2]{FIG/logo-mba-2016.png}} % número de página a la derecha

\renewcommand{\headrulewidth}{0.4pt} % grosor de la línea de la cabecera
\renewcommand{\footrulewidth}{0.4pt} % grosor de la línea del pie

\begin{document}



\pagebreak
\vspace*{4mm} 

\begin{center}
\begin{huge}
\textbf{\\Econometría: Ejercicios 1}
\end{Large}\\
%{\large 2023-2}\\
\end{center}



\section{Regresión Lineal Simple}
\textbf{Ejercicio 5:} Se tienen los siguientes datos de salario, años de educación y experiencia laboral de 5 individuos:
\begin{center}
\begin{tabular}{|c|c|c|c|}
\hline
\textbf{Individuo} & \textbf{Educ} ($X_1$) & \textbf{Exp} ($X_2$)  & \textbf{Salario ($Y$)}\\
\hline
1 & 12 & 14 & 1000 \\
2 & 14 & 15 & 1200 \\
3 & 15 & 17 & 1300 \\
4 & 16 & 15 & 1200 \\
5 & 18 & 16 & 1500 \\
\hline
\end{tabular}
\end{center}

\begin{enumerate}
    \item Planteando un modelo donde el salario sólo dependa de la educación se tiene la siguiente ecuación:
    $$\widehat{Salario} = \beta_0 + \beta_1 \cdot Educ$$
    Calcule los valores de los estimadores $\hat{\beta_0}$ y $\hat{\beta_1}$ y SST, SSE, SSR, $R^2$.
    \\
    
    \textbf{solución}


\textcolor{red}{

\[
\bar X=\frac{12+14+15+16+18}{5}=15,\qquad
\bar Y=\frac{1000+1200+1300+1200+1500}{5}=1240.
\]

\begin{center}
\begin{tabular}{|r|r|r|r|r|r|r|}
\hline
$X_i$ & $Y_i$ & $X_i-\bar X$ & $Y_i-\bar Y$ & $(X_i-\bar X)^2$ & $(Y_i-\bar Y)^2$ & $(X_i-\bar X)(Y_i-\bar Y)$\\ \hline
12 & 1000 & $-3$ & $-240$ & $9$  & $57{,}600$ & $720$ \\
14 & 1200 & $-1$ & $-40$  & $1$  & $1{,}600$  & $40$  \\
15 & 1300 & $0$  & $60$   & $0$  & $3{,}600$  & $0$   \\
16 & 1200 & $1$  & $-40$  & $1$  & $1{,}600$  & $-40$ \\
18 & 1500 & $3$  & $260$  & $9$  & $67{,}600$ & $780$ \\ \hline
\multicolumn{2}{|r|}{\textbf{Totales}} & & & $\mathbf{20}$ & $\mathbf{132{,}000}$ & $\mathbf{1{,}500}$\\ \hline
\end{tabular}
\end{center}

\[
S_{xx}=20,\quad S_{xy}=1{,}500.
\]


\textbf{Estimadores MCO:}
\[
\hat\beta_1=\frac{S_{xy}}{S_{xx}}=\frac{1500}{20}=75,\qquad
\hat\beta_0=\bar Y-\hat\beta_1\bar X=1240-75\cdot 15=115.
\]
\[
\widehat{Y}_i=115+75\,X_i.
\]

\textbf{Descomposición de suma de cuadrados y } 

\begin{center}
\begin{tabular}{|r|r|r|r|r|}
\hline
$Educ$ & $Y$ & $\widehat{Y}=115+75\,Educ$ & $e=Y-\widehat{Y}$ & $e^{2}$ \\ \hline
12 & 1000 & 1015 & $-15$  & $225$ \\
14 & 1200 & 1165 & $35$   & $1225$ \\
15 & 1300 & 1240 & $60$   & $3600$ \\
16 & 1200 & 1315 & $-115$ & $13225$ \\
18 & 1500 & 1465 & $35$   & $1225$ \\ \hline
\multicolumn{4}{|r|}{\textbf{SSR} $=\sum e_i^2$} & \textbf{19500} \\ \hline
\end{tabular}
\end{center}



$R^2$:
\[
SST=132{,}000,\qquad SSR=19{,}500
\]
\[
SSE=SST-SSR=132{,}000-19{,}500=12{,}500,\qquad
R^2=\frac{SSR}{SST}=\frac{112{,}500}{132{,}000}\approx 0.8523.
\]

\textbf{Lectura:} $\widehat{Salario}=115+75\cdot Educ$. Un año adicional de educación se asocia, en promedio, con \$75 más de salario. El ajuste explica $\approx 85.23\%$ de la variación de $Y$.}


\item En una empresa se realizó un curso de un software importante para el desarrollo del trabajo, siendo este obligatorio con selección aleatoria para algunos trabajadores, y tambien los no escogidos aleatoriamente podían inscribirse de forma voluntaria.
Sea $Y$ el salario, $D\in\{0,1\}$ el tratamiento (tomó el curso) y $R\in\{0,1\}$ la asignación aleatoria.
Se observan 8 individuos (4 tratados, 4 no tratados). Entre los tratados, 3 son aleatorios ($R=1$); entre los no tratados, 2 son aleatorios.

\begin{center}
\begin{tabular}{|c| c |c |c|}
\hline
id & $Y$ & $D$ & $R$\\
\hline
1 & 950 & 1 & 1\\
2 & 900 & 1 & 1\\
3 & 880 & 1 & 1\\
4 & 820 & 1 & 0\\
5 & 780 & 0 & 1\\
6 & 740 & 0 & 1\\
7 & 700 & 0 & 0\\
8 & 650 & 0 & 0\\
\hline
\end{tabular}
\end{center}

\begin{enumerate}
\item Calcule la diferencia ingenua $\hat\beta_{\text{ing}}=\overline{Y}_{D=1}-\overline{Y}_{D=0}$.
\\
\textbf{Solución}
\\
El beta ingenuo (sesgado, antes de considerar aleatoriedad) esta dado por:
$$\beta_{ing}=E[Y|D=1]-E[Y|D=0]$$
$$E[Y|D=1] = \frac{950+900+880+820}{4}=887.5$$
$$E[Y|D=0]=\frac{780+740+700+650}{4}=717.5$$
$$\beta_{ing}=887.5-717.5=170$$

\item En la submuestra $R=1$, calcule $\hat\beta_{R} = \overline{Y}_{D=1,R=1}-\overline{Y}_{D=0,R=1}$.
\\
\textbf{Solución}
\\
Ahora considerando la aleatoriedad, consideraremos solo los casos que fueron escogidos aleatoriamente y calcularemos el $\beta$.
$$\hat\beta = E[{Y}|{D=1,R=1}]-E[{Y}|{D=0,R=1}]$$
$$ E[{Y}|{D=1,R=1}] = \frac{950+900+880}{3}=910$$
$$ E[{Y}|{D=0,R=1}] = \frac{780+740}{2}=760$$
$$\hat{\beta}= 910-760=150$$
\item Interprete $\hat\beta$ como estimador cercano al $ATE$ (y al $ATT$ con efectos homogéneos). Calcule el sesgo de selección en este caso.
\\
\textbf{Solución}
\\
Como el calculo de $\hat{\beta}$ es sobre aleatorios, podemos decir que $\hat{\beta}=ATE=ATT$, y si comparamos con el valor ingenuo, podemos encontrar el valor del sesgo:
$$Sesgo=\hat{\beta_{ing}}-\hat{\beta}=170-150=20$$
El sesgo era igual a 20, y estaba sobreestimando el valor real de $\beta$.
\item Discuta por qué $\hat\beta_{\text{ing}}$ puede estar sesgado cuando $D$ no es aleatorio.
\\
\textbf{Solución}
\\
La decisión de participar de un curso, cuando es de forma voluntaria, depende de varios factores, como inteligencia del individuo, motivación, entre otros... Y estas características podrían influir tambien en tener un salario mayor.
\end{enumerate}

\item Queremos realizar la siguiente regresión de $Y=\beta_0+\beta_1X+\epsilon$, pero solo podemos tener una versión aproximada de $Y^*=Y+c$, con $c$ una constante desconocida, y $X^*=X+u$, siendo $u$ error clásico ($cov(X,u)=0, var(u)=\sigma^2_u$), por lo que sólo podemos tener una versión imprecisa de la regresión, $Y^* = \tilde\beta_0+\tilde\beta_1X^*+v$. Desmuestra que 
$$\hat{\tilde\beta}_1=\hat{\beta_1}\frac{Var(X)}{Var(X)+\sigma^2_u}$$
\\
\textbf{Solución}
\\
necesitamos encontrar 
$$\tilde{\hat{\beta}}_1=\frac{Cov(X^*,Y^*)}{Var(X^*)}$$
sabemos que $Y^*=Y+c$, con $c$ una constante desconocida, y $X^*=X+u$, siendo $u$ error clásico. Entonces
$$\tilde{\hat{\beta}}_1=\frac{Cov(X+u,Y+c)}{Var(X+u)}$$
$$\tilde{\hat{\beta}}_1=\frac{Cov(X,Y)+Cov(X,c)+Cov(Y,u)+ Cov(c,u))}{Var(X)+Var(u)}$$
Sabemos que la varianza de una constante con cualquier valor es 0, y además como $u$ es error clásico, no correlaciona con $Y$, por lo que queda
$$\tilde{\hat{\beta}}_1=\frac{Cov(X,Y))}{Var(X)+\sigma^2_u}$$
Luego sabemos que $Y=\beta_0+\beta_1X+\epsilon$, por lo que
$$\tilde{\hat{\beta}}_1=\frac{Cov(X,\beta_0+\beta_1X+\epsilon))}{Var(X)+\sigma^2_u}$$
distribuyendo:
$$\tilde{\hat{\beta}}_1=\frac{Cov(X,\beta_0)+Cov(X,\beta_1X)+Cov(X,\epsilon))}{Var(X)+\sigma^2_u}$$
$Cov(X,\beta_0)=0$ porque $\beta_0$ es constante y $Cov(X,\epsilon)=0$, por el supuesto de $E[\epsilon|X]=0$, entonces:
$$\tilde{\hat{\beta}}_1=\frac{Cov(X,\beta_1X)}{Var(X)+\sigma^2_u}=\beta_1\frac{Var(X)}{Var(X)+\sigma^2_u}$$

\end{enumerate}


\section{Formulario}
\begin{itemize}

\item  Estimadores MCO para RLS.
{
$$\hat{\beta}_0 = \bar{y} - \hat{\beta}_1 \bar{x} $$
$$\hat{\beta}_1 = \frac{\sum_{i=1}^{n}(x_i - \bar{x})(y_i-\bar{y})}{\sum_{i=1}^{n}(x_i - \bar{x})^2}=\frac{Cov(x,y)}{Var(x)}$$}

\item \[ R^2  = 1 - \frac{SSR}{SST} \] 
\item  $$SST = \sum_{i=1}^{n} (y_i - \bar{y})^2 \quad \quad SSR = \sum_{i=1}^{n}  \hat{u}_i^2 $$
    \item \[ \bar{R}^2 =  1 - \frac{(n-1)SSR}{(n-k-1)SST} \]

\end{itemize}

\begin{figure}
    \centering
    \includegraphics[width=0.8\linewidth]{Figuras/normal_critic.png}
\end{figure}





\end{document}